%Chargement des paquets
\usepackage{amsmath}
\usepackage{amsfonts}
\usepackage{amssymb}
\usepackage{enumerate}
\usepackage{mathtools}
\usepackage{bbm}
\usepackage{xparse, etoolbox}
\usepackage{enumerate}
\usepackage{enumitem}
\usepackage{mathabx}
\usepackage{babel}[french]

%environment
\usepackage{amsthm}
\usepackage{keytheorems}
\usepackage{hyperref} 

%Interface théorème
\newkeytheorem{lem}[
	name = Lemme
]

\newkeytheorem{prop}[
	name = Proposition
]

\newkeytheorem{thm}[
	name = Théorème
]

\newkeytheorem{coro}[
	name = Corollaire
]

\renewcommand*{\proofname}{Démonstration}

\theoremstyle{definition}
\newtheorem{defn}{Définition}

\theoremstyle{definition}
\newtheorem*{nota}{Notation}

\theoremstyle{definition}
\newtheorem*{limi}{Liminaire}

\theoremstyle{remark}
\newtheorem{rmq}{Remarque}

\theoremstyle{plain}
\newtheorem*{fait}{Fait}


%Ensembles de nombres
\newcommand{\N} {\mathbb{N}}
\newcommand{\Ne}{\N^\ast}
\newcommand{\Z} {\mathbb{Z}}
\newcommand{\Q} {\mathbb{Q}}
\newcommand{\R} {\mathbb{R}}
\newcommand{\Rb}{\overline{\mathbb{R}}}
\newcommand{\Rp}{\R_+}
\newcommand{\Rm}{\R_-}
\newcommand{\K} {\mathbb{K}}
\newcommand{\Ket} {\mathbb{K^{\ast}}}
\newcommand{\Cx}{\mathbb{C}}

%Ensembles, tribus
\newcommand{\A}{\mathbb{A}}
\renewcommand{\P}{\mathcal{P}}
\newcommand{\T}{\mathcal{T}}
\newcommand{\F}{\mathcal{F}}
\newcommand{\C} {\mathcal{C}}
\newcommand{\ouv}{\mathcal{O}}
\newcommand{\B}{\mathcal{B}}
\newcommand{\M}{\mathcal{M}}
\newcommand{\etag}{\mathcal{E}}
\newcommand{\etagOp}{\etag^{+}(\Omega)}
\newcommand{\etagO}{\etag(\Omega)}
%%Lp avant quotientage
\newcommand{\Lic}{\mathcal{L}^1}
\newcommand{\Lpc}{\mathcal{L}^p}
\newcommand{\Linfc}{\mathcal{L}^{\infty}}
\newcommand{\Lifull}{\Lic(\Omega, \B, m)}
\newcommand{\Lpfull}{\Lpc(\Omega, \B, m)}
\newcommand{\Linffull}{\Linfc(\Omega, \B, m)}
%%Lp après quotientage
\newcommand{\Li}{L^1}
\newcommand{\Ld}{L^2}
\newcommand{\Lp}{L^p}
\newcommand{\Linf}{L^{\infty}}
\newcommand{\Listd}{\Li(\Omega, m)} 
\newcommand{\Ldstd}{\Ld(\Omega, m)} 
\newcommand{\Lpstd}{\Lp(\Omega, m)} 

%Quelques opérateurs
\DeclareMathOperator{\codim}{codim}
\DeclareMathOperator{\rg}{rg}
\DeclareMathOperator{\tr}{tr}
\DeclareMathOperator{\vect}{Vect}
\DeclareMathOperator{\ima}{Im}
\DeclareMathOperator{\id}{Id}
\DeclareMathOperator{\sign}{sign}
\DeclareMathOperator{\ext}{ext}
\newcommand{\ind}{\mathbbm{1}}
\newcommand*{\spr}[2]{\left(\,#1\,\middle|\,#2\,\right)}
\newcommand{\interior}[1]{%
  {\kern0pt#1}^{\mathrm{o}}%
}
\DeclareMathOperator{\card}{card}
\DeclareMathOperator{\ppcm}{ppcm}

%Commandes de pure flemme
\newcommand{\veps}{\varepsilon}
\newcommand{\intab}{\int_{a}^{b}}
\newcommand{\intR}{\int_{-\infty}^{\infty}}
\newcommand{\intinf}{\int_{- \infty}^{+ \infty}}
\newcommand{\equi}{\Leftrightarrow}
\newcommand{\Kev}{$\K$-espace vectoriel }
\newcommand{\Kevs}{$\K$-espaces vectoriels }
\newcommand{\Rev}{$\R$-espace vectoriel }
\newcommand{\Revs}{$\R$-espaces vectoriels }
\newcommand{\Cev}{$\Cx$-espace vectoriel }
\newcommand{\Cevs}{$\Cx$-espaces vectoriels }
\newcommand{\evn}{(E, \norm{.})}
\newcommand{\interf}[2]{[#1,#2]}
\newcommand{\limb}{\lim\limits_}
\newcommand{\lebn}{\lambda_n}
\newcommand{\pp}{\text{~presque partout}}
\newcommand{\derivp}[2]{\frac{\partial #1}{\partial #2}}
\newcommand{\famiu}{(u_i)_{i \in I}}
\newcommand{\sev}{\text{~s.e.v.}}
\newcommand{\Mnp}{M_{n,p}(\K)}
\newcommand{\Mn}{M_{n}(\K)}
\newcommand{\tq}{\text{~t.q.~}}
\newcommand{\mq}{~montrer~que~}
\newcommand{\et}{\quad \text{et} \quad}
\newcommand{\ou}{\text{~ou~}}
\newcommand{\Kx}{\K[X]}


%Commandes de gars chelou
\renewcommand{\subset}{\subseteq}

%Commande restriction, ty duckduckgo
\def\restr#1#2{\mathchoice
              {\setbox1\hbox{${\displaystyle #1}_{\scriptstyle #2}$}
              \restraux{#1}{#2}}
              {\setbox1\hbox{${\textstyle #1}_{\scriptstyle #2}$}
              \restraux{#1}{#2}}
              {\setbox1\hbox{${\scriptstyle #1}_{\scriptscriptstyle #2}$}
              \restraux{#1}{#2}}
              {\setbox1\hbox{${\scriptscriptstyle #1}_{\scriptscriptstyle #2}$}
              \restraux{#1}{#2}}}
\def\restraux#1#2{{#1\,\smash{\vrule height .8\ht1 depth .85\dp1}}_{\,#2}} 

%Quelques normes
\DeclarePairedDelimiter\abs{\lvert}{\rvert}
\DeclarePairedDelimiter\labs{\bigl\lvert}{\bigr\rvert}
\DeclarePairedDelimiter\norm{\lVert}{\rVert}
\DeclarePairedDelimiterXPP\normi[1]{}\lVert\rVert{_1}{\ifblank{#1}{\:\cdot\:}{#1}}
\DeclarePairedDelimiterXPP\normd[1]{}\lVert\rVert{_2}{\ifblank{#1}{\:\cdot\:}{#1}}
\DeclarePairedDelimiterXPP\normp[1]{}\lVert\rVert{_p}{\ifblank{#1}{\:\cdot\:}{#1}}
\DeclarePairedDelimiterXPP\normq[1]{}\lVert\rVert{_q}{\ifblank{#1}{\:\cdot\:}{#1}}
\DeclarePairedDelimiterXPP\norminf[1]{}\lVert\rVert{_\infty}{\ifblank{#1}{\:\cdot\:}{#1}}

%Bracket en angle
\DeclarePairedDelimiter\angl{\langle}{\rangle}

%Set, parenthèses
\newcommand{\set}[1]{\{ #1 \}}
\newcommand{\bigset}[1]{\bigl\{ #1 \bigr\}}
\newcommand{\Bigset}[1]{\Bigl\{ #1 \Bigr\}}
\newcommand{\bigpar}[1]{\bigl( #1 \bigr)}
\newcommand{\Bigpar}[1]{\Bigl( #1 \Bigr)}

%Opitmisation des opérateurs de comparaison
\DeclarePairedDelimiter\tio{o (}{)}
\DeclarePairedDelimiter\gro{O (}{)}


\pagestyle{fancy}

\fancyhead[C]{Matrice et déterminant de Gram}
\fancyhead[R]{}

\begin{document}

\underline{Source :} Gourdon - Algèbre - page 274 et Grifone - Algèbre linéaire - page 268 (cas euclidien) \\

\underline{Leçons possibles :} 
\begin{itemize}
\item 149 : déterminant. Exemples et applications.
\item 157 : matrices symétriques réelles, matrices hermitiennes.
\item 206 : exemples d'utilisation de la notion de dimension finie en analyse.
\end{itemize}

\begin{nota}
Dans ce développement, $E$ est un espace préhilbertien (réel ou complexe), son produit scalaire (éventuellement hermitien) 
est désigné par $\angl{ \cdot, \cdot}$. 
Cet espace est muni de la norme $\norm{.}$ induite par son produit scalaire. \\
Pour toute matrice $M$ à coefficients complexes, on note $M^{\ast}$ sa matrice transconjugée. 
Une matrice égale à sa transconjugée est dite hermitienne.
\end{nota}

\begin{defn}[Matrice de Gram]
Soient $x_1, \dots, x_n$ des vecteurs de $E, n \in \Ne$. 
On appelle matrice de Gram de $x_1, \dots, x_n$ la matrice carrée dont le coefficient d'indice $(i,j)$ est $\angl{x_i, x_j}$. 
Le déterminant de Gram est le déterminant de cette matrice.
\end{defn}

\begin{nota}
Pour une famille $x_1, \dots, x_n$, on note usuellement $G(x_1, \dots, x_n)$ le déterminant de Gram associé.
\end{nota}

\begin{thm}
Toute matrice de Gram est hermitienne positive. Réciproquement, toute matrice hermitienne positive est une matrice de Gram.
De plus, la matrice de Gram de $n$ vecteurs $x_1, \dots, x_n$ est définie si, et seulement si, la famille de vecteurs est libre.
\end{thm}

\begin{proof}

\begin{enumerate}[label=(\alph*)]

\item
Soient $x_1, \dots, x_n$ des vecteurs de $E$ et $G$ leur matrice de Gram. Soit $F = \vect(x_1, \dots, x_n)$ et $m = \dim{F}$. 
On fixe $\B$ une base orthonormée de $F$ (toujours possible en dimension finie), et on note $X_i$ le vecteur colonne des coordonnées
de $x_i$ dans la base $\B$ et $M$ la matrice $(m \times n)$ dont les colonnes sont les $X_i$, de sorte que :
\[ \angl{x_i , x_j} = X_i^{\ast} X_j \et G = M^{\ast}M . \]
On en déduit que la matrice $G$ est hermitienne ($G^{\ast} = M ^{\ast} M^{\ast\ast} = M^{\ast} M = G$) et positive (pour tout vecteur 
colonne $X$, on a $X^{\ast}GX = X^{\ast} M^{\ast} M X = (MX)^{\ast} MX = \norm{MX}$). 

\item
Réciproquement, si $A$ est une matricienne hermitienne positive, d'après le théorème spectral, il existe une matrice unitaire $U$
telle que $D = U^{\ast} A U$ soit diagonale réelle. De plus, comme $A$ est positive, les coefficients de $D$ sont positifs 
(pour s'en convaincre, on remarquera que le coefficient $i$ de $D$ est le résultat du produit $C_i^{\ast} A C_i$ où $C_i$ est la colonne
$i$ de la matrice $U$). On peut donc écrire $D= (D')^2$ avec $D'$ une autre matrice diagonale à coefficients réels positifs, et donc :
\[ A = (U D' U^{\ast} ) (U D' U^{\ast}) = B^2 = B^{\ast} B , \]
car $B$ est une matrice hermitienne. Donc, si on désigne par $X_i$ les colonnes de la matrice $B$, le coefficient $(i,j)$ de $A$ est 
\[ (X_i)^{\ast} X_j = \angl{x_i, x_j} \]
et $A$ est une matrice de Gram. 

\item
On reprend les notations du point (a) pour étudier le cas défini, on note donc $G = M^{\ast}M$. 
La matrice $G$ est définie si, et seulement $(\forall Y \in M_{n,1}(\Cx), Y^{\ast} G Y = 0 \Rightarrow X = 0)$ ce qui est aussi équivalent à
$(\forall Y \in M_{n,1}(\Cx), \norm{M Y} = 0 \Rightarrow Y = 0)$ ce qui est une formulation matricielle de la liberté de la famille 
$x_1, \dots, x_n$.

\end{enumerate}
\end{proof}

L'intérêt des matrices de Gram est, que leurs déterminants permettent de calculer la distance d'un point à un s.e.v. en utilisant le 
théorème suivant.

\begin{thm}
Soit $E$ un espace préhilbertien, $V$ un sous-espace de $E$ muni d'une base quelconque $(e_1, \dots, e_n)$. 
Alors, la distance de $x$ à $V$ définie comme $d = \inf_{v \in V} \norm{x-v}$ vérifie :
\[ d^2 = \dfrac{G(e_1, \dots, e_n, x)}{G(e_1, \dots, e_n)} . \]
\end{thm}

\begin{proof}
On suppose acquis le fait que $d$ vérifie $d = \norm{x - p_V(x)}$ où $p_V(x)$ est la projection orthogonale de $x$ sur $V$. 
On note $y = p_V(x)$ et on a :
\[ \forall i \in \set{1, \dots, n}, \quad \angl{e_i, x } = \angl{e_i, y} \et \norm{x} = \norm{x-y} + \norm{y} , \]
car $\angl{ x-y, y} = 0$. On en déduit que :
\[ \left | \begin{array}{c | c}
\begin{matrix}
\angl{e_1, e_1} & \dots & \angl{e_1, e_n} \\
\vdots & & \vdots \\
\angl{e_n, e_1} & \dots & \angl{e_n, e_n}
\end{matrix} & \begin{matrix} \angl{e_1, x} \\ \vdots \\ \angl{e_n, x} \end{matrix} \\
\hline 
\begin{matrix} \angl{x, e_1} & \dots & \angl{x, e_n} \end{matrix} & \angl{x, x}
\end{array} \right | 
= \left | \begin{array}{c | c}
\begin{matrix}
\angl{e_1, e_1} & \dots & \angl{e_1, e_n} \\
\vdots & & \vdots \\
\angl{e_n, e_1} & \dots & \angl{e_n, e_n}
\end{matrix} & \begin{matrix} \angl{e_1, y} \\ \vdots \\ \angl{e_n, y} \end{matrix} \\
\hline 
\begin{matrix} \angl{y, e_1} & \dots & \angl{y, e_n} \end{matrix} & \angl{y,y} + d^2
\end{array} \right | \]
Soit, en utilisant la linéarité du déterminant par rapport à la dernière colonne :
\[ G(e_1, \dots, e_n, x) =
\left | \begin{array}{c | c}
\begin{matrix}
\angl{e_1, e_1} & \dots & \angl{e_1, e_n} \\
\vdots & & \vdots \\
\angl{e_n, e_1} & \dots & \angl{e_n, e_n}
\end{matrix} & \begin{matrix} \angl{e_1, y} \\ \vdots \\ \angl{e_n, y} \end{matrix} \\
\hline 
\begin{matrix} \angl{y, e_1} & \dots & \angl{y, e_n} \end{matrix} & \angl{y,y}
\end{array} \right |
+ \left | \begin{array}{c | c}
\begin{matrix}
\angl{e_1, e_1} & \dots & \angl{e_1, e_n} \\
\vdots & & \vdots \\
\angl{e_n, e_1} & \dots & \angl{e_n, e_n}
\end{matrix} & \begin{matrix} 0 \\ \vdots \\ 0 \end{matrix} \\
\hline 
\begin{matrix} \angl{y, e_1} & \dots & \angl{y, e_n} \end{matrix} & d^2
\end{array} \right | \]
Or, le premier déterminant de cette somme est nul car $y \in \vect(e_1, \dots, e_n)$, d'où il vient :
\[ G(e_1, \dots, e_n, x) = d^2 G(e_1, \dots, e_n) , \]
en développant le second déterminant de la somme selon la dernière colonne.
\end{proof}

\begin{rmq}
Il convient de noter que le premier théorème nous assure que $G(e_1, \dots, e_n)$ est non nul dans le second théorème 
(car la famille de vecteurs est une base, en particulier, est libre).
\end{rmq}

\end{document}